\scanpage{27}
\mspage{8}

\[
\rho,f,g,u,v,w\qquad\in M
\]
de $M$, satisfaisant les relations ($**$), ou encore tels qu'il existe $p,q,r,s\in\mathbb N$ avec
\[
\boxed{\varphi^p(\rho gfu)=\varphi^q(\rho gv)=\varphi^r(\rho w)=\varphi^s\psi}
\]

Ceci posé, on doit poser la relation
\[
\boxed{[w,gv][v,fu][w,gfu]^{-1}=1},
\]

Pour l'expliciter, considérons de façon générale
\diagram{\begin{tikzcd}[ampersand replacement=\&]e\arrow[r,shift left=.55ex,"\lambda"]\arrow[r,shift right=.55ex,"\mu"']\&\xi=a\sigma\end{tikzcd}}
donc $e=a\sigma\lambda=a\sigma\mu$ i.e.
\[
\boxed{\sigma\lambda\sim\sigma\mu\sim\psi},
\]

et choisissons $\lambda',\mu'\in M$ tels que $\lambda\lambda'=\mu\mu'=\unsure{\pi|\beta}$ de sorte qu'on a un diagramme commutatif
\diagram{\begin{tikzcd}[ampersand replacement=\&]\&e\arrow[dr,"\lambda"]\&\\ \eta\arrow[ur,"\lambda'"]\arrow[dr,"\mu'"']\&\&\xi\\ \&e\arrow[ur,"\mu"']\&\end{tikzcd}}
\[
\begin{aligned}\text{avec}\quad\eta&=e\lambda'=\xi\lambda\lambda'=a\sigma\lambda\lambda'\deleted{\unsure{=\pi}}\inserted{=a\sigma\pi}\\ \cdot&=e\mu'=\xi\mu\mu'=a\sigma\mu\mu'=a\sigma\pi\end{aligned}
\]

on va \deleted{choisir} prendre une flèche \diagram{\begin{tikzcd}[ampersand replacement=\&]e\arrow[r,"\nu"]\&\eta\end{tikzcd}} donc un $\nu\in M$ tel que
% letterhead
\par\smallskip
{\footnotesize\noindent\begin{minipage}[t]{.46\linewidth}\raggedright
MINISTERE DE L'EDUCATION NATIONALE\\ ACADEMIE D'AIX-MARSEILLE\\ \rule{2.6em}{.3pt}\\ Inspection Académique\\ de Vaucluse\\ \rule{2.6em}{.3pt}\\ Tél. : 90 82 99 76\\ Référence : Secrétariat n\textsuperscript{o}
\end{minipage}\hfill%
\begin{minipage}[t]{.51\linewidth}\raggedright
AVIGNON, le\\ L'Inspecteur d'académie,\\ directeur des services départementaux\\ de l'Education Nationale de Vaucluse,\\ \deleted{à M}
\end{minipage}\par}

% archive sheet number: 26
